% BHU PYQ -- Manually transcribed & error-corrected
% Paper: BCH-224 : Business Mathematics
% Exam: B.Com. (Hons.) Semester IV Examination, 2023-24

\documentclass[11pt,a4paper]{article}
\usepackage[a4paper,top=2.5cm,bottom=2.5cm,left=2.8cm,right=2.8cm,headheight=14pt]{geometry}
\usepackage{amsmath,amssymb}
\usepackage[T1]{fontenc}
\usepackage[utf8]{inputenc}
\usepackage{lmodern,microtype,fancyhdr,enumitem,hyperref,lastpage,parskip}

\hypersetup{
    colorlinks=true,
    urlcolor=black,
    linkcolor=black,
    pdftitle={BCH-224: Business Mathematics -- BHU B.Com. (Hons.) Sem IV 2023-24}
}

\pagestyle{fancy}
\fancyhf{}
\renewcommand{\headrulewidth}{0.4pt}
\renewcommand{\footrulewidth}{0.4pt}
\fancyhead[L]{\small   \textit{Banaras Hindu University}}
\fancyhead[R]{\small   \textit{BCH-224: Business Mathematics}}
\fancyfoot[C]{\small Page  \thepage\ of \pageref{LastPage}}


\newlist{parts}{enumerate}{2}
\setlist[parts,1]{label=(\roman*),leftmargin=*,itemsep=5pt,topsep=3pt}
\setlist[parts,2]{label=(\alph*),leftmargin=*,itemsep=3pt,topsep=2pt}

\newcommand{\pts}[1]{\hfill   \textit{[#1]}}


\begin{document}


\begin{center}
  {\large\bfseries BANARAS HINDU UNIVERSITY}\\[2pt]
  {
\normalsize B.Com. (Hons.) --- Semester IV Examination, 2023-24}\\[6pt]
  \rule{0.8\linewidth}{0.5pt}\\[6pt]
  {\large\bfseries Paper No. BCH-224: Business Mathematics}\\[4pt]
  {
\normalsize Time: 3 Hours\hfill Maximum Marks: 70}
\end{center}

\rule{\linewidth}{0.4pt}

\smallskip



\noindent   \textit{Instructions: Attempt all questions. All questions carry equal marks (14 marks each).}

\bigskip




\noindent   \textbf{Question 1.}

\begin{parts}
  \item The $p \text{th}$ term of an A.P. is $q$ and $q \text{th}$ term is $p$, prove that $m \text{th}$ term will be $p+q-m$.\pts{7}
  \item Sum of 8 terms of an A.P. is 64 and sum of 16 terms is 256. Find the sum of $n$ terms.\pts{7}
\end{parts}

\centerline{   \textbf{OR}}


\begin{parts}
  \item 7th and 13th terms of an A.P. are 13 and 25 respectively. Find the sum of 10 terms.\pts{7}
  \item Sum of 3 terms of an A.P. is 15 and their product is 105. Find the numbers.\pts{7}
\end{parts}

\medskip\rule{\linewidth}{0.2pt}




\noindent   \textbf{Question 2.}

\begin{parts}
  \item If the sum of $n$ terms of a G.P. $\frac{2}{9} - \frac{1}{3} + \frac{1}{2} - \dots$ is $-\frac{1261}{576}$, find $n$.\pts{7}
  \item A person desires to save Rs. 2 in January, Rs. 4 in February, Rs. 8 in March, Rs. 16 in April and so on upto the end of the year. Determine:
  
\begin{parts}
    \item How much amount he will save during one year?\pts{3.5}
    \item What will he save in the month of October?\pts{3.5}
  \end{parts}
\end{parts}

\centerline{   \textbf{OR}}


\begin{parts}
  \item $(p+q) \text{th}$ term of a G.P. is $m$ and $(p-q) \text{th}$ term is $n$. Find $p \text{th}$ term.\pts{7}
  \item Arithmetic mean between two numbers is 10 and their geometric mean is 8. Find the numbers.\pts{7}
\end{parts}

\medskip\rule{\linewidth}{0.2pt}




\noindent   \textbf{Question 3.}

\begin{parts}
  \item If ${}^{2n+1}P_{n-1} : {}^{2n-1}P_n = 3:5$, find $n$.\pts{7}
  \item A box contains 7 red, 6 white and 4 blue balls. How many selections of 3 balls can be made so that:
  
\begin{parts}
    \item None is red.\pts{3.5}
    \item One is of each colour.\pts{3.5}
  \end{parts}
\end{parts}

\centerline{   \textbf{OR}}


\begin{parts}
  \item Prove that: ${}^{n-1}P_r + r \cdot {}^{n-1}P_{r-1} = {}^nP_r$.\pts{7}
  \item If ${}^{2n}C_3 : {}^nC_2 = 44:3$, find $n$.\pts{7}
\end{parts}

\medskip\rule{\linewidth}{0.2pt}




\noindent   \textbf{Question 4.}

\begin{parts}
  \item Find the term independent of $x$ in the expansion of $\left(x^2 + \frac{1}{x}\right)^{12}$.\pts{7}
  \item If in the expansion of $(1+x)^{24}$ the coefficient of $r \text{th}$ term and $(r+4) \text{th}$ term are equal, find $r$.\pts{7}
\end{parts}

\centerline{   \textbf{OR}}

Prove that the middle term in the expansion of $(1+x)^{2n}$ is $\frac{1 \cdot 3 \cdot 5 \dots (2n-1)}{n!} 2^n x^n$.\pts{14}

\medskip\rule{\linewidth}{0.2pt}




\noindent   \textbf{Question 5.}

\begin{parts}
  \item If $f(x) = \frac{x-1}{x+1}$, then find the value of $\frac{f(x) - f(y)}{1 + f(x)f(y)}$.\pts{4}
  \item Evaluate: $\lim_{x  o \infty} \frac{3x^2 + 5x - 5}{5x^2 + 6x + 2}$.\pts{4}
  \item Differentiate w.r.t. $x$: $\frac{(x+1)^2 \sqrt{x-1}}{(x+4)^3 e^x}$.\pts{6}
\end{parts}

\centerline{   \textbf{OR}}


\begin{parts}
  \item Simplify: $\int \frac{\sqrt{x} + 3\sqrt{x}}{5x\sqrt{x}} \, dx$.\pts{4}
  \item Simplify: $\int \frac{1}{\sqrt{x+3} - \sqrt{x}} \, dx$.\pts{4}
  \item Simplify: $\int \frac{(x+1)e^x}{(x+2)^2} \, dx$.\pts{6}
\end{parts}

\vfill
\rule{\linewidth}{0.4pt}

\begin{center}\small
\href{https://bhu-exam.vercel.app/}{Click here to visit our website}\\[4pt]
\href{https://me-aryan.vercel.app/}{Click here to visit our contributor}\\[4pt]
\href{https://quashan-me.vercel.app/}{Click here to visit our contributor}
\end{center}

\end{document}
